\section{Execution Cycle \& Memory Handling}
The interaction between the deterministic orchestration layer and the generative model is governed by a stateful execution cycle. Unlike standard single-turn LLM interactions, the Galois framework maintains a persistence layer across iterations to maximize recall and prevent redundancy. This logic is encapsulated within the \textbf{GaloisExecutor} class and its specialized internal memory component, \textbf{GaloisMemory}.

\subsection{The $\texttt{GaloisMemory}$ Management Class}

\subsubsection{Purpose and Architecture}
The $\texttt{GaloisMemory}$ class serves as the dedicated state management component for the physical execution strategies ($\texttt{Table-Scan}$ and $\texttt{Key-Scan}$). Its primary responsibilities are:
\begin{enumerate}
    \item \textbf{History Accumulation:} Storing the unique records (tuples) extracted by the LLM chain across multiple iterative steps.
    \item \textbf{Deduplication Enforcement:} Ensuring that only unique records are preserved in the memory structure ($\mathcal{M}$) to maintain data integrity and prevent redundant processing.
    \item \textbf{Context Provision:} Formatting the accumulated records ($\mathcal{M}$) into a structured history ($\mathcal{H}$) for injection into subsequent LLM prompts, guiding the model toward novel data discovery.
\end{enumerate}
The class inherits from $\texttt{BaseMemory}$ and $\texttt{BaseModel}$, structuring the internal state via Pydantic fields.

\subsubsection{Internal State Variables}
The memory state is maintained by the following core attributes:
\begin{itemize}
    \item $\texttt{memory}$ (\texttt{List[Dict[str, Any]]}): The primary storage structure containing the complete collection of unique records (tuples) retrieved thus far.
    \item $\texttt{key\_values}$ (\texttt{set}): A hash set storing the unique key value combinations (as Python tuples) for rapid, $O(1)$ deduplication lookups. This structure is central to enforcing uniqueness across iterations.
    \item $\texttt{tokens}$ (\texttt{int}): An aggregate counter tracking the total number of tokens processed by the LLM across all invocations for the current extraction task.
    \item $\texttt{time}$ (\texttt{float}): An aggregate measure of the total processing time (latency) associated with the LLM interactions.
\end{itemize}

\subsubsection{Core Methods and Functionality}

\begin{enumerate}
    \item \textbf{$\texttt{load\_memory\_variables(inputs)}$}
    \begin{itemize}
        \item \textit{Function:} Prepares the accumulated records for use as historical context in the next LLM prompt.
        \item \textit{Implementation:} Converts the internal $\texttt{memory}$ list into a compact JSON string, which is returned under the key $\texttt{"history"}$. This serialized format is robust for prompt injection.
    \end{itemize}

    \item \textbf{$\texttt{save\_context(inputs, outputs)}$}
    \begin{itemize}
        \item \textit{Function:} The central method for processing new records generated by an LLM invocation and updating the memory state.
        \item \textit{Procedure:} Iterates through the $\texttt{outputs["response"]}$ (newly generated records). For each record, it extracts the unique key values using $\texttt{\_get\_key\_values}$. If the key is not present in $\texttt{self.key\_values}$, the record is appended to $\texttt{self.memory}$, and the key is added to $\texttt{self.key\_values}$.
        \item \textit{Convergence Check:} Compares the size of $\texttt{self.key\_values}$ before and after processing the new records. If the set size remains unchanged, it indicates that \textit{no new unique tuples} were found, triggering a $\texttt{NoNewTuplesFound}$ exception to terminate the iterative process (per Algorithms \ref{alg:table_scan} and \ref{alg:key_scan}).
        \item \textit{Metric Update:} Aggregates the $\texttt{tokens}$ and $\texttt{time}$ metrics from the $\texttt{outputs}$ into the class state.
    \end{itemize}

    \item \textbf{$\texttt{clear()}$}
    \begin{itemize}
        \item \textit{Function:} Resets the memory state, typically called after an extraction process completes.
        \item \textit{Implementation:} Empties the $\texttt{self.memory}$ list and $\texttt{self.key\_values}$ set, resetting $\texttt{tokens}$ and $\texttt{time}$ to zero.
    \end{itemize}

    \item \textbf{$\texttt{\_get\_key\_values(record, keys)}$}
    \begin{itemize}
        \item \textit{Function:} An internal helper method to extract the composite key values from a single record dictionary.
        \item \textit{Implementation:} Takes a $\texttt{record}$ (dict) and a list of $\texttt{keys}$ (column names) and returns a $\texttt{tuple}$ containing the corresponding values. This tuple is used as the hashable item for deduplication checks in the $\texttt{key\_values}$ set.
    \end{itemize}
\end{enumerate}


\subsection{The Iterative Execution Loop and Termination Criteria}

Both the $\texttt{Table-Scan}$ and $\texttt{Key-Scan}$ physical operators rely on a fundamental Iterative Execution Loop (lines 4-19 in Algorithms \ref{alg:table_scan} and \ref{alg:key_scan}) to overcome the LLM context window limitations and maximize data extraction coverage. The core mechanism involves repeatedly querying the LLM and dynamically updating the prompt with historical data.

\subsubsection{Loop Mechanism: Pagination \& History Injection}
In both strategies, the iterative process is driven by the dynamic insertion of previously acquired data into the subsequent prompt. This mechanism ensures two critical functions:
\begin{enumerate}
    \item \textbf{Deduplication:} The LLM is implicitly guided to avoid regenerating records already present in the injected history ($\mathcal{H}$).
    \item \textbf{Novelty Search:} The context directs the LLM to explore uncollected portions of the solution space, thereby maximizing the cardinality of the final result set ($\mathcal{R}_{final}$).
    \item To enhance retrieval efficiency and prevent the LLM from wandering randomly, the iterative mechanism implements a \textbf{value-based pagination} strategy: In addition to the history the system identifies the last retrieved value (\texttt{last\_val}) from the memory.
    The prompt logic explicitly instructs the LLM to generate unique tuples that strictly follow \texttt{last\_val} (alphabetically or numerically). This transforms the extraction into a pseudo-sequential scan, improving coverage compared to simple history injection.
\end{enumerate}
The difference lies in the content of the history: $\texttt{Table-Scan}$ injects a history of complete tuples, while $\texttt{Key-Scan}$ injects only the unique key values ($\mathcal{K}$) collected so far.

\subsection{Termination Criteria (Convergence Conditions)}
The iterative loop is governed by two distinct termination criteria, ensuring both safety against infinite execution and functional completeness (convergence). The loop terminates if either condition is met:

\begin{enumerate}
    \item \textbf{Maximum Iteration Limit ($I_{max}$):}
    \begin{itemize}
        \item \textit{Condition:} The loop counter ($i$) reaches the maximum permissible value defined by the input parameter $I_{max}$, which is configurable via the system's YAML configuration file or directly supplied by the user.
        \item \textit{Purpose:} This serves as a fail-safe mechanism to bound the token consumption and execution time, preventing excessive resource usage in cases where the LLM continually produces marginally unique or redundant data.
    \end{itemize}

    \item \textbf{Data Convergence ($\mathcal{R}_{new} = \emptyset$ or $\mathcal{K}_{new} = \emptyset$):}
    \begin{itemize}
        \item \textit{Condition:} After an LLM invocation and subsequent parsing/deduplication step, the system determines that \textbf{no new unique records (tuples)} were added to the memory (for $\texttt{Table-Scan}$), or \textbf{no new unique key values} were collected (for $\texttt{Key-Scan}$).
        \item \textit{Mechanism:} This condition is detected by the $\texttt{save\_context}$ function of the $\texttt{GaloisMemory}$ class, typically by comparing the cardinality of the $\texttt{key\_values}$ set before and after processing the LLM's latest output.
        \item \textit{Purpose:} This is the primary functional exit criterion, indicating that the LLM has likely exhausted its capacity to discover novel data relevant to the query $Q$, signaling functional convergence and allowing for premature loop termination for efficiency.
    \end{itemize}
\end{enumerate}
The system employs the $\texttt{break}$ statement within the $\texttt{While}$ loop when convergence is detected, overriding the $I_{max}$ constraint if full coverage is achieved earlier.